
Time series data from digital platforms exhibit temporal dynamics that may follow distinct trajectory patterns, yet systematic methods for characterizing these patterns remain underexplored. This work proposes a two-level hierarchical methodology combining PELT (Pruned Exact Linear Time) changepoint detection with DTW (Dynamic Time Warping) based clustering for trajectory analysis. The first level identifies discrete phase transitions using PELT with BIC penalty selection, while the second level groups trajectories by shape similarity using DTW-based TimeSeriesKMeans clustering. Experiments on 500 time series from HuggingFace demonstrate that trajectories partition into clusters with silhouette score 0.352 and bootstrap Jaccard stability 0.82. PELT changepoint detection identifies at least one statistically significant changepoint in 81\% of time series. Shape descriptor analysis reveals that three of four descriptors (growth ratio, changepoint count, derivative variance) differentiate clusters with variance ratios exceeding 2.0, while peak timing does not discriminate between clusters. When mapping the four empirical clusters to five proposed behavioral archetypes, only two archetypes (slow\_burn and revival) are recovered, with multiple clusters mapping to the same archetype. These results validate the two-level hierarchical methodology for trajectory analysis while indicating that the empirical cluster structure is simpler than the hypothesized five-archetype taxonomy.
